\documentclass[12pt]{article}
\usepackage[margin=0.6in,top=0.65in,bottom=0.5in]{geometry}
\usepackage{lmodern}
\usepackage{microtype}
\usepackage{fancyhdr}
\usepackage{titlesec}
\usepackage{enumitem}
\usepackage{lastpage}
\usepackage[hidelinks]{hyperref}

\setlength{\parindent}{0pt}
\setlength{\parskip}{0.22em}
\setlength{\headheight}{26pt}
\raggedbottom
\titleformat{\section}{\normalsize\bfseries\scshape}{}{0pt}{}
\titleformat{\subsection}{\normalsize\bfseries}{}{0pt}{}
\titlespacing*{\section}{0pt}{0.4em}{0.15em}
\titlespacing*{\subsection}{0pt}{0.3em}{0.1em}
\pagestyle{fancy}
\fancyhf{}
\fancyhead[C]{\footnotesize Lit Example Reader \quad Assignment ID: U1L8LE \quad Page \thepage\ of \pageref{LastPage}}
\renewcommand{\headrulewidth}{0.5pt}

\newenvironment{playpassage}{\begin{quote}\small\setlength{\parskip}{0.12em}\setlength{\parindent}{0pt}}{\end{quote}}
\newcommand{\speaker}[1]{\textbf{#1.}}

\begin{document}

\begin{center}
{\LARGE\bfseries Week 8 Lit Example Reader}\par\vspace{0.02em}
{\large\scshape Macbeth: Fair, Foul, and Shifting Security}\par\vspace{0.08em}
\rule{0.78\textwidth}{0.5pt}
\end{center}

\textbf{Purpose:} Macbeth passages for Week 8: expose term shifts that make prophecy and persuasion look like stable proof.

\textbf{What to watch for:} Repeated words that do not keep one idea. Ask which sense is operating, and whether an inference still holds when the sense is fixed.

\section*{Before the Passages: The Week 8 Logic Tool}

\begin{itemize}[leftmargin=1.3em,itemsep=0.08em,topsep=0.1em,parsep=0.04em]
\item \textbf{Equivocation:} one word, two ideas, one argument.
\item \textbf{Ambiguous middle:} a linking term that only looks shared.
\item \textbf{Fix the term:} state both senses; rewrite; retest the conclusion.
\end{itemize}

\section*{Passage 1: Act 1, Scene 1 and Scene 3 --- Fair, Foul, and Greeting}

\textbf{Context:} The witches open the play by joining opposite words. Later they greet Macbeth with titles that sound like facts and promises.

\begin{playpassage}
\speaker{All} Fair is foul, and foul is fair:\\
Hover through the fog and filthy air.\\
\speaker{First Witch} All hail, Macbeth! hail to thee, Thane of Glamis!\\
\speaker{Second Witch} All hail, Macbeth! hail to thee, Thane of Cawdor!\\
\speaker{Third Witch} All hail, Macbeth, that shalt be king hereafter!\\
\speaker{Banquo} Good sir, why do you start, and seem to fear\\
Things that do sound so fair?
\end{playpassage}

\subsection*{Reader Notes}

Here is your most assisted example. \textit{Fair} is doing unstable work. In the witches' chant, fair and foul are forced together so moral and sensory language blur. In Banquo's line, \textit{fair} can mean attractive, promising, or fortunate-sounding. The titles also mix present fact (Glamis) with future claim (king hereafter). Week 8 asks whether a later inference from these words keeps one sense of \textit{fair}, \textit{hail}, or promised greatness---or whether Macbeth treats a shifting verbal surface as if it were one solid middle term.

\textbf{Reading task:} Circle every use of \textit{fair}. Mark which titles are present fact and which are future claim. In the margin, write two possible senses of \textit{fair} in this passage.

\newpage

\section*{Passage 2: Act 4, Scene 1 --- Security in the Apparitions}

\textbf{Context:} The apparitions give Macbeth statements that sound like safety. Watch the key words that seem to guarantee him.

\begin{playpassage}
\speaker{Second Apparition} Be bloody, bold, and resolute; laugh to scorn\\
The power of man, for none of woman born\\
Shall harm Macbeth.\\
\speaker{Third Apparition} Macbeth shall never vanquish'd be until\\
Great Birnam wood to high Dunsinane hill\\
Shall come against him.\\
\speaker{Macbeth} That will never be:\\
Who can impress the forest, bid the tree\\
Unfix his earth-bound root? Sweet bodements! good!\\
Rebellion's head, rise never till the wood\\
Of Birnam rise, and our high-plac'd Macbeth\\
Shall live the lease of nature, pay his breath\\
To time and mortal custom.
\end{playpassage}

\subsection*{Reader Notes}

This passage shifts more work to you. Macbeth hears words of security and treats them as a stable guarantee. Mark the phrases that sound impossible. Decide what ordinary sense Macbeth gives \textit{none of woman born} and \textit{Birnam wood}, and what different sense later events use. Do not complete the whole worksheet diagnosis here---only mark the pressure words and possible sense split.

\textbf{Reading task:} Underline the two security promises. Beside each, write one sense that would make Macbeth feel safe. Save the full fixed-meaning rewrite for the worksheet.

\section*{How to Use These Passages on the Worksheet}

The worksheet asks you to treat prophetic language as an argument with terms. Keep senses separate. A repeated word is not automatically a stable middle.

\end{document}
